% © 2026 Ing. David Jaroš — CC BY-NC-ND 4.0
%
% This work is licensed under a Creative Commons Attribution-NonCommercial-NoDerivatives
% 4.0 International License (CC BY-NC-ND 4.0).
%
% License History: Earlier drafts (up to v0.3) were released under CC BY 4.0.
% From v0.4 onward, all material is released under CC BY-NC-ND 4.0 to protect
% the integrity of the theoretical work during ongoing academic development.
%
% See LICENSE.md for full license text.
%
% File: canonical/alpha/gamma_entropy_alpha_refinement_status.tex
% Purpose: Status document for the Gamma/prime-factorization entropy interpolation
%          audit of the alpha effective potential. Records numerical findings,
%          mathematical identities, and the bracket property of V1 and V3
%          relative to the measured fine-structure constant.
% Date: 2026-05-11
% Related: canonical/alpha/prime_factorization_entropy_potential.tex
%          canonical/alpha/ALPHA_MASTER_STATUS.md
%          tools/scan_gamma_entropy_lambda.py
%          reports/gamma_entropy_alpha_interpolation_audit.md

% UBT-AI-PROVENANCE-BEGIN
% schema: ubt-ai-provenance/v1
% tier: B_machine_verified
% ai_assistance: disclosed
% human_review: machine-verification
% editorial_responsibility: Ing. David Jaroš
% policy: ../../AI_PROVENANCE.md
% notice: Machine-verified against named sources or verifiers; individual attestation is not claimed.
% UBT-AI-PROVENANCE-END

\ifdefined\INCLUDEMODE
\else
  \documentclass[12pt,a4paper]{article}
  \usepackage[a4paper, margin=2.5cm]{geometry}
  \usepackage{amsmath, amssymb, amsthm}
  \usepackage{mathtools}
  \usepackage{hyperref}
  \usepackage{xcolor}
  \usepackage{booktabs}
  \usepackage{longtable}
  \usepackage{mdframed}
  \usepackage{tcolorbox}
  \tcbuselibrary{skins,breakable}
  \usepackage{array}

  \newtheorem{theorem}{Theorem}[section]
  \newtheorem{lemma}[theorem]{Lemma}
  \newtheorem{proposition}[theorem]{Proposition}
  \newtheorem{corollary}[theorem]{Corollary}
  \theoremstyle{definition}
  \newtheorem{definition}[theorem]{Definition}
  \theoremstyle{remark}
  \newtheorem{remark}[theorem]{Remark}

  \newmdenv[backgroundcolor=yellow!20, linecolor=orange,          linewidth=1.5pt]{gapbox}
  \newmdenv[backgroundcolor=green!15,  linecolor=green!60!black,  linewidth=1.5pt]{resultbox}
  \newmdenv[backgroundcolor=red!8,     linecolor=red!60,          linewidth=1.5pt]{warnbox}
  \newmdenv[backgroundcolor=blue!6,    linecolor=blue!40,         linewidth=1.5pt]{infobox}
  \newmdenv[backgroundcolor=gray!8,    linecolor=gray!60,         linewidth=1pt  ]{statusbox}
  \newmdenv[backgroundcolor=orange!10, linecolor=orange!70,       linewidth=1pt  ]{obsbox}

  % Status labels
  \newcommand{\STDLZ}{\textcolor{gray!70!black}{\textbf{[STD/L0]}}}
  \newcommand{\INTERP}{\textcolor{blue!60!black}{\textbf{[INTERP]}}}
  \newcommand{\OBS}{\textcolor{orange!70!black}{\textbf{[OBS]}}}
  \newcommand{\DERIVCAND}{\textcolor{orange!80!black}{\textbf{[DERIVATION CANDIDATE]}}}
  \newcommand{\OPEN}{\textcolor{orange}{\textbf{[OPEN]}}}
  \newcommand{\PROVED}{\textcolor{green!60!black}{\textbf{[PROVED]}}}
  \newcommand{\PHENOM}{\textcolor{purple!70!black}{\textbf{[PHENOM]}}}
  \newcommand{\NUMONLY}{\textcolor{red!50!black}{\textbf{[NUMERIC\_ONLY]}}}
  \newcommand{\COND}{\textcolor{blue}{\textbf{[COND]}}}
  \newcommand{\Lone}{\textbf{[L1]}}
  \newcommand{\Lzero}{\textbf{[L0]}}

  \newcommand{\BRam}{B_{\mathrm{Ram}}}
  \newcommand{\Veff}{V_{\mathrm{eff}}}
  \newcommand{\Neff}{N_{\mathrm{eff}}}

  \title{\textbf{Gamma/Prime-Factorization Entropy Refinement:\\
         Alpha Potential Interpolation Status}\\
         \large Audit of the $V_\lambda$ Bracket Property and $\lambda_{\mathrm{fit}}$}
  \author{Ing.\ David Jaro\v{s}}
  \date{2026-05-11}

  \begin{document}
  \maketitle
  \tableofcontents
  \newpage
\fi

%──────────────────────────────────────────────────────────────────────────────
\section{Status Discipline}
\label{sec:status}
%──────────────────────────────────────────────────────────────────────────────

\begin{warnbox}
\textbf{Mandatory status declarations:}
\begin{itemize}
  \item Alpha ($\alpha$) is \textbf{NOT DERIVED}.
  \item $\BRam$ is \OBS{} — a numerical observation, not derived from $S[\Theta]$.
  \item Gap G137-B is \OPEN{} and is not closed by this analysis.
  \item $\lambda_{\mathrm{fit}}$ is a fit parameter; it is \OBS{} unless derived.
  \item All identities from standard mathematics are labeled \STDLZ{}.
  \item The prime-factorization entropy interpretation is labeled \INTERP{}.
\end{itemize}
\end{warnbox}

This document records the status of the Gamma/prime-factorization entropy
interpolation audit for the UBT alpha potential.  It does not claim any new
derivation of $\alpha$ or of $B$.  All numerical results come from
\texttt{tools/scan\_gamma\_entropy\_lambda.py} (mpmath, 80 decimal places).

%──────────────────────────────────────────────────────────────────────────────
\section{Potential Definitions}
\label{sec:potentials}
%──────────────────────────────────────────────────────────────────────────────

\subsection{V1 — Original potential}

\begin{equation}
  V_1(n) = n^2 - B \cdot n\log n,
  \qquad
  V_1'(n) = 2n - B(\log n + 1) = 0.
  \tag{V1, \Lone}
\end{equation}

\subsection{V3 — Gamma-refined potential}

\begin{equation}
  V_3(n) = n^2 - B\cdot\bigl(\log\Gamma(n+1) + n\bigr),
  \qquad
  V_3'(n) = 2n - B(\psi(n+1) + 1) = 0,
  \tag{V3, \DERIVCAND}
\end{equation}
where $\psi = \psi^{(0)}$ is the digamma function.

\subsection{V\_lambda — Interpolated potential}

\begin{equation}
  E_\lambda(n) = (1-\lambda)\cdot n\log n + \lambda\cdot(\log\Gamma(n+1) + n),
  \label{eq:Elambda}
\end{equation}
\begin{equation}
  V_\lambda(n) = n^2 - \BRam \cdot E_\lambda(n).
  \label{eq:Vlambda}
\end{equation}

The stationarity condition $V_\lambda'(n) = 0$ is:
\begin{equation}
  2n = \BRam\,\bigl[(1-\lambda)(\log n + 1) + \lambda(\psi(n+1) + 1)\bigr].
  \label{eq:stat_lambda}
\end{equation}

\subsection{Reference value $\BRam$ \OBS}

\begin{obsbox}
$\BRam = 12^{3/2} \cdot 2^{1/8} \cdot \vartheta_3(0\mid i)^{1/4}$
\quad \OBS{} — not derived from $S[\Theta]$.
\end{obsbox}

Numerically (80 decimal places):
\begin{equation}
  \BRam = 46.280872383180703498930750708\ldots
  \label{eq:BRam_value}
\end{equation}

%──────────────────────────────────────────────────────────────────────────────
\section{Mathematical Identities}
\label{sec:identities}
%──────────────────────────────────────────────────────────────────────────────

\begin{resultbox}
\textbf{Identity 1} \STDLZ{}.  For every $n \in \mathbb{N}$:
\begin{equation}
  \log\Gamma(n+1) = \log(n!).
\end{equation}

\textbf{Identity 2 — Legendre's formula} \STDLZ{}.  For every $n \in \mathbb{N}$:
\begin{equation}
  \log(n!) = \sum_{\substack{p \text{ prime},\, m \geq 1 \\ p^m \leq n}}
             \left\lfloor \frac{n}{p^m} \right\rfloor \log p.
  \label{eq:legendre}
\end{equation}

\textbf{Identity 3 — Stirling expansion} \STDLZ{}.
\begin{equation}
  \log\Gamma(n+1) + n = n\log n + \tfrac{1}{2}\log(2\pi n) + \tfrac{1}{12n} + O(n^{-3}).
  \label{eq:stirling_refined}
\end{equation}

\textbf{Identity 4 — Digamma derivative} \STDLZ{}.
\begin{equation}
  \frac{d}{dn}\bigl[\log\Gamma(n+1) + n\bigr] = \psi(n+1) + 1,
  \qquad
  \psi(n+1) = \log n + \tfrac{1}{2n} - \tfrac{1}{12n^2} + O(n^{-4}).
  \label{eq:digamma_asymp}
\end{equation}
\end{resultbox}

%──────────────────────────────────────────────────────────────────────────────
\section{High-Precision Numerical Results}
\label{sec:numerics}
%──────────────────────────────────────────────────────────────────────────────

All values computed at \texttt{mpmath.mp.dps = 80}.

\begin{center}
\begin{tabular}{@{}lll@{}}
\toprule
Quantity & Value & Status \\
\midrule
$\vartheta_3(0\mid i)$ & $1.08643481121330801457531612151\ldots$ & \STDLZ \\
$|\eta(i)|$ & $0.768225422326056659002594179576\ldots$ & \STDLZ \\
$\BRam$ & $46.280872383180703498930750708\ldots$ & \OBS \\
$n_1$ (V1 stationary) & $136.9890996341081484047897\ldots$ & \Lone \\
$n_3$ (V3 stationary) & $137.0905214130678236021181\ldots$ & \DERIVCAND \\
$n_3 - n_1$ & $0.10142177895967520232826\ldots$ & computed \\
$\alpha_{\mathrm{exp}}^{-1}$ (CODATA 2018) & $\approx 137.036$ & \PHENOM \\
\bottomrule
\end{tabular}
\end{center}

\subsection{Bracket property}

\begin{resultbox}
\textbf{Bracket property} (numeric, \OBS{} since depends on $\BRam$):
\begin{equation}
  n_1 < \alpha_{\mathrm{exp}}^{-1} < n_3,
  \qquad
  n_1 = 136.989\ldots,\quad
  \alpha_{\mathrm{exp}}^{-1} = 137.036\ldots,\quad
  n_3 = 137.091\ldots
  \label{eq:bracket}
\end{equation}
The measured fine-structure constant inverse lies between the V1 and V3
stationary points.  The physical $\alpha^{-1}$ lies approximately
\textbf{46\% of the way} from $n_1$ to $n_3$.
\end{resultbox}

\subsection{Lambda fit values}

Two definitions of $\lambda_{\mathrm{fit}}$ are used; they differ by $\approx 0.0002$
due to the nonlinearity of $n^*(\lambda)$:

\paragraph{Definition A (exact stationarity).}
The unique $\lambda$ such that $V_\lambda'(n) = 0$ at $n = \alpha_{\mathrm{exp}}^{-1}$:
\begin{equation}
  \lambda_{\mathrm{fit}}^{\mathrm{exact}} =
    \frac{2\alpha_{\mathrm{exp}}^{-1}/\BRam - (\log\alpha_{\mathrm{exp}}^{-1}+1)}
         {(\psi(\alpha_{\mathrm{exp}}^{-1}+1)+1) - (\log\alpha_{\mathrm{exp}}^{-1}+1)}
  = 0.46221754271946\ldots
  \quad \OBS
  \label{eq:lambda_exact}
\end{equation}

\paragraph{Definition B (linear fractional position).}
\begin{equation}
  \lambda_{\mathrm{fit}}^{\mathrm{frac}} =
    \frac{\alpha_{\mathrm{exp}}^{-1} - n_1}{n_3 - n_1}
  = 0.46241990993372\ldots
  \quad \OBS
  \label{eq:lambda_frac}
\end{equation}

Both values are \OBS{} and have no derivation from $S[\Theta]$.

%──────────────────────────────────────────────────────────────────────────────
\section{Information-Theoretic Interpretation \INTERP}
\label{sec:interp}
%──────────────────────────────────────────────────────────────────────────────

\begin{infobox}
\textbf{Entropy interpretation} \INTERP{}.  By Legendre's formula
(equation~\eqref{eq:legendre}), $\log(n!) = \log\Gamma(n+1)$ is the
\emph{total logarithmic information content} of integers $1,\ldots,n$ when each
integer is encoded as a composite symbol over the prime alphabet with weights
$w_p = \log p$.

The Stirling leading term $n\log n$ (identity~\eqref{eq:stirling_refined})
is the leading approximation to this entropy.  The interpolation $E_\lambda(n)$
therefore interpolates between:
\begin{itemize}
  \item $\lambda = 0$: the Stirling leading-term approximation $n\log n$,
  \item $\lambda = 1$: the exact prime-factorization entropy $\log(n!) + n$.
\end{itemize}
The measured $\alpha^{-1}$ lies at $\lambda \approx 0.462$ of the way between
these two structural limits.  The sub-percent difference between $n_1$ and $n_3$
(and hence between the two entropy models) is produced by the Stirling subleading
term $\tfrac{1}{2}\log(2\pi n) \approx 3.38$ at $n = 137.036$.
\end{infobox}

This interpretation is \INTERP{} only: it does not derive $B$, $\lambda$, or
$\alpha$ from $S[\Theta]$.

%──────────────────────────────────────────────────────────────────────────────
\section{Candidate Constant Comparison}
\label{sec:candidates}
%──────────────────────────────────────────────────────────────────────────────

The table below compares $\lambda_{\mathrm{fit}}^{\mathrm{exact}} = 0.46221754\ldots$
against candidate constants.  Status ``\NUMONLY'' means: numerically close but
no UBT derivation known.

\begin{center}
\begin{tabular}{@{}lllll@{}}
\toprule
Constant & Value & $|\text{err}|$ & rel.\% & Status \\
\midrule
$37/80$ & $0.46250$ & $2.8\times10^{-4}$ & $0.061\%$ & \NUMONLY \\
$6/13$ & $0.46154$ & $6.8\times10^{-4}$ & $0.147\%$ & \NUMONLY \\
$\tfrac{1}{2} - \tfrac{1}{12\pi}$ & $0.47347$ & $1.1\times10^{-2}$ & $2.4\%$ & \NUMONLY \\
$3/(2\pi)$ & $0.47746$ & $1.5\times10^{-2}$ & $3.3\%$ & \NUMONLY \\
$1/2$ & $0.50000$ & $3.8\times10^{-2}$ & $8.2\%$ & \NUMONLY \\
$\sqrt{2}-1$ & $0.41421$ & $4.8\times10^{-2}$ & $10.4\%$ & \NUMONLY \\
$1/e$ & $0.36788$ & $9.4\times10^{-2}$ & $20.4\%$ & \NUMONLY \\
(all modular) & $\eta(i), \vartheta_3,\ldots$ & $>0.15$ & $>30\%$ & \NUMONLY \\
\bottomrule
\end{tabular}
\end{center}

\begin{warnbox}
\textbf{Key finding}: No candidate constant has a derivation from UBT.
The closest match, $37/80 = 0.4625$, is within $0.06\%$ of
$\lambda_{\mathrm{fit}}^{\mathrm{exact}}$, but this coincidence has no
structural interpretation and must be classified \NUMONLY{}.
\end{warnbox}

%──────────────────────────────────────────────────────────────────────────────
\section{Gap G137-B Status}
\label{sec:gap}
%──────────────────────────────────────────────────────────────────────────────

\begin{gapbox}
\textbf{Gap G137-B} \OPEN{}.

The following remain unresolved:
\begin{enumerate}
  \item Derivation of $\BRam \approx 46.28$ from $S[\Theta]$ without using $\alpha$.
  \item Derivation of the interpolation parameter $\lambda_{\mathrm{fit}} \approx 0.462$
        from any UBT mechanism (RG, Stirling, modular form, determinant regularization).
  \item Identification of the physical mechanism selecting the observed $\alpha^{-1}$
        within the bracket $[n_1, n_3]$.
\end{enumerate}

The bracket property~\eqref{eq:bracket} is a \OBS{} numerical finding, not a proof.
The physical fine-structure constant lying between the asymptotic entropy model
($\lambda = 0$) and the Gamma-refined entropy model ($\lambda = 1$) is structurally
suggestive but does not close G137-B.
\end{gapbox}

%──────────────────────────────────────────────────────────────────────────────
\section{Conclusion}
\label{sec:conclusion}
%──────────────────────────────────────────────────────────────────────────────

\begin{statusbox}
\textbf{Overall verdict}: \textbf{PLAUSIBLE\_OBS}

The physical alpha inverse lies between the stationary points of the asymptotic
entropy potential (V1, $\lambda=0$) and the exact Gamma-entropy potential
(V3, $\lambda=1$), with:
\[
  n_1 = 136.989099\ldots < \alpha_{\mathrm{exp}}^{-1} = 137.035999\ldots < n_3 = 137.090521\ldots
\]
The Stirling identities and Legendre formula provide exact mathematical grounding
for this structure (\STDLZ).  The entropy interpretation (\INTERP) explains why
$n\log n$ is a natural quantity.  Neither $\BRam$ nor $\lambda_{\mathrm{fit}}$ is
derived from $S[\Theta]$.

\textbf{Alpha is NOT DERIVED.  Gap G137-B is OPEN.}
\end{statusbox}

%──────────────────────────────────────────────────────────────────────────────
\ifdefined\INCLUDEMODE
\else
  \end{document}
\fi
